\documentclass[11pt]{article}
\usepackage{amsmath,amssymb,amsthm,mathtools,hyperref}

\begin{document}
\section*{ecDNA generation and function}
\label{prob:249}
\label{rank:106}
\par\smallskip

\noindent\textbf{Domain.} Nonchromosomal inheritance inference \quad \textbf{Retained status/source context.} ecDNA-associated amplification and inheritance have substantial established results, including coordinated inheritance in 2024. Origins and context-dependent quantitative impact remain an active programme.
\paragraph{Definitions and assumptions.} Let $N_t=(N_{1t},\ldots,N_{mt})\in\mathbb{N}_0^m$ be copy counts of $m$ predeclared ecDNA species in a cell generation. A finite-dimensional parameter specifies inheritance/formation kernel $K_\theta$ on this full countable space, with $\sum_{n'}K_\theta(n'\mid n,c)=1$. Expansion weights $w_\eta(n,c)>0$ are bounded. Existing lineage/snapshot data observe counts through a stated error model and an explicit censored overflow category above detection limits; censoring does not truncate the biological kernel. Formation, segregation and selection are distinct. Independent binomial segregation across species is a model option, not a default fact.
\paragraph{Problem statement.} Determine which features of $K$ and $w$ can be separated by available data. A population-distribution update is
    \[P_{t+1}(n')=\frac{\sum_nP_t(n)w_\eta(n,c)K_\theta(n'\mid n,c)}
    {\sum_nP_t(n)w_\eta(n,c)}.\]
    Identify the parameter equivalence classes giving the same observed copy-number distributions, and bound correlations attributable to inheritance rather than post-division selection. The target is retrospective statistical explanation.
\paragraph{Additional binding assumptions and resolution criterion.} Derive $w(n,c)K(n'\mid n,c)$ from the expected offspring counting measure. For such a representation, let $L$ be the number of offspring counted over the declared interval and let $X_1,\ldots,X_L$ be their states at the declared counting phase. Under the declared conditional offspring law, write
\[
M(n'\mid n,c)=\mathbb{E}\!\left[\sum_{j=1}^{L}\mathbf{1}_{\{X_j=n'\}}\,\middle|\,n,c\right],\qquad
w(n,c)=\sum_{n'}M(n'\mid n,c),\qquad
K(n'\mid n,c)=\frac{M(n'\mid n,c)}{w(n,c)}.
\]
This is the offspring-weighted kernel. A randomly selected daughter's transition law and a parent's expected offspring number need not factor this way if offspring number and daughter state are dependent. The interval, counting phase and inclusion of death or selection in $M$ must be declared; the weighted kernel is not automatically segregation before selection. This representation does not select whether $P_t$ denotes normalized expected counts, a deterministic population limit or random finite-population frequencies. That choice remains unspecified. Normalizing expected counts differs in general from averaging normalized random frequencies; a lineage-sampling law also differs from population sampling. Any finite-population interpretation must address extinction and any survival conditioning, and any population-limit interpretation must supply its limit and supporting assumptions.

The stated strict positivity of $w$ excludes zero-expected-offspring types in an expected-offspring representation; it does not exclude stochastic population extinction. The zero count vector describes an ecDNA-free cell, not an extinct cell population, and need not be absorbing when formation is allowed. For a probability law $P_t$, bounded positive weights and a stochastic kernel already give a finite positive denominator in the displayed update. Tail and moment control must instead be specified for the requested readouts and censoring bounds, over the declared horizon and parameter domain, with propagation and quantitative tail bounds where needed; no particular tail family is imposed here. Censoring all high counts cannot by itself distinguish arbitrary tail-dependent formation or selection rates. Equivalence classes must retain any admissible common positive rescaling of $w$ that leaves the declared observable laws unchanged; normalized copy-number distributions alone do not determine absolute reproductive output. Absolute growth or offspring observations may supply additional information, but their availability is not specified.

\paragraph{Instance specifications and remaining choices.} As required by the binding catalogue conventions, an instance must declare the available dataset or admissible observation designs, candidate families and parameter domains, initial laws, observation times and horizon, error model and overflow thresholds. Observational equivalence means equality of the full observable laws under the declared designs, not equality of one sample likelihood or merely of snapshot marginals when lineage data are included. The instance must also declare the operational identity and detection rules for ecDNA species, the biological population, the phase of copy counting relative to replication and division, and the generation interval. Here $c$ denotes context whose domain, observation status and evolution remain to be specified; no fixed, time-varying or inherited convention is selected. Whether counts and context form a closed Markov state, or require an augmented-state model or an explicitly justified approximation, remains an instance requirement.

Both requested deliverables remain binding: characterize observational equivalence and bound the declared inheritance-versus-selection correlation target. For the latter, declare the random variables, correlation functional, conditioning variables, reference population, measurement phases, zero-variance convention and statistical attribution rule. A sister-correlation target requires a joint offspring-state law, which a single-offspring kernel does not determine. The source material does not select among within-cell, sister or mother--daughter correlations, nor specify the inheritance-versus-selection comparison. These choices remain unresolved; no causal interpretation follows from statistical attribution alone.

\paragraph{Alternative formulations.} (i) Direct mother--daughter observations constrain inheritance more strongly than snapshots. (ii) Single-species and multispecies kernels differ because co-segregation can preserve correlation before selection. (iii) ecDNA count and causal tumour-fitness contribution are separate targets; association with expansion does not itself identify a causal function.
\paragraph{Obstructions and partial progress.} Turner et al. establish ecDNA-associated heterogeneity; Hung et al. demonstrate coordinated multispecies inheritance and modeling, updating older assumptions of independent random assortment. The active challenge therefore does not mean inheritance principles are wholly unmodeled. Formation and selection can produce similar snapshot distributions, and copy-number-dependent detection adds another ambiguity. A useful result would characterize that non-identifiability and test explicitly restricted kernels against existing lineage evidence.
\paragraph{Sources (supplied provenance; not all rechecked).}
\begin{itemize}
\item Cancer Grand Challenges, active challenge index, accessed 11 October 2026. Research-agenda evidence, not a mathematical conjecture. \url{https://www.cancergrandchallenges.org/challenges/active-challenges}
\item Cancer Grand Challenges, Extrachromosomal DNA; active programme, description set in 2020. \url{https://www.cancergrandchallenges.org/challenges/active-challenges/extrachromosomal-dna}
\item K. M. Turner et al., Extrachromosomal oncogene amplification drives tumour evolution and genetic heterogeneity (2017); established amplification/heterogeneity role. \url{https://www.nature.com/articles/nature21356}
\item K. L. Hung et al., Coordinated inheritance of extrachromosomal DNAs in cancer cells (2024); inheritance and quantitative-model progress. \url{https://www.nature.com/articles/s41586-024-07861-8}
\end{itemize}


\section*{Binding catalogue conventions}
This edition revises both supplied catalogues and preserves original problem numbers 1--330. It adds one globally unique label \verb|\label{rank:x}| per problem, where smaller $x$ means greater provisional difficulty. A rank label is an anchor at that problem's section: \verb|\ref{rank:x}| returns the section number, not the literal difficulty number. The visible \texttt{rank:x} text and the index give the difficulty number itself.

The entries are \emph{parameterized mathematical research schemes}, including empirical prediction, causal identification, historical reconstruction, theory and engineering. They are not 330 certified unsolved mathematical conjectures. A rigorous instance fixes the objects and restrictions below; changing those objects changes the question. Established identities, finite calculations, non-identification constructions and solved special cases are not being posed anew as open. Some residual tasks concern informative bounds, biological realization or external validation rather than a new theorem.

Every main definition and problem statement received a mathematical editorial review. Each entry now has an additional binding-conditions paragraph; direct corrections are separately recorded in the accompanying audit. The alternative formulations inherit the applicable conventions, but are different targets unless equivalence is proved. Where an old broad phrase conflicts with an explicit restriction below it, the explicit restriction controls.

\textbf{Literature boundary.} Thirteen selective primary-literature checks cover eighteen entries in this revision; their access levels and scope are recorded explicitly. The supplied bibliography and its earlier source-access descriptions are retained as provenance, not represented as newly verified. Neither this review nor a source's description proves the absence of later solutions. The date of this edition is 13 October 2026; there is no claim of exhaustive current-status certification for all 330 entries.

\subsection*{Binding conventions for a concrete instance}
\paragraph{Objects, quantifiers and non-vacuity.}
An instance declares measurable state, input, output and parameter spaces; a fixed candidate family; admissible initial laws, interventions and environments; a finite horizon or specified asymptotic regime; observation and sampling laws; and performance metrics, positive normalization scales and positive tolerances. These are chosen before fitting or testing, not adjusted after seeing failures. Shared parameters must genuinely be shared across the declared validation domain. Finite dimension alone does not restrict arbitrary real-number encodings: specify a finite library, bounded description length, or regularity/complexity bounds. Exclude constant or no-signal solutions through a fixed informative-output or baseline-improvement requirement where appropriate. A universally quantified result cannot be established by choosing one trivial favorable instance.

\paragraph{Probability laws and data.}
Every expectation or probability refers to the specified stochastic law, including initial uncertainty, process noise, observation noise and sampling as appropriate. A deterministic equation is not a probability model without a declared random input or initial law. ``Observed'' targets in population-level formulas mean population observables $g_j^*$, not noisy estimates asserted to be exact. Let $\widetilde g_j$ be independent validation estimates and suppose
\[
\mathbb{P}\{\lvert\widetilde g_j-g_j^*\rvert\le r_j\text{ for every }j\}\ge1-\alpha.
\]
Then an uncertainty-aware sufficient certification rule for positive target tolerances $\varepsilon_j$ is
\[
\max_j\frac{|\widehat g_j-\widetilde g_j|+r_j}{\varepsilon_j}\le1.
\]
On that simultaneous confidence event it implies every population error is at most its tolerance, by the triangle inequality. Vector/norm versions use corresponding valid simultaneous bounds. Training uncertainty must also be included when the claim averages over repeated training data. Evaluation on finitely many conditions does not establish a continuous-domain supremum without covering and regularity bounds.

\paragraph{Dynamics and well-posedness.}
For ODEs and SDEs give coefficient domains, initial laws and existence/uniqueness conditions in the intended solution sense, or keep the entire admitted solution set. For stochastic jump models specify nonnegative off-diagonal rates, diagonal normalization, rate orientation and non-explosion. For PDEs specify spatial domains, boundary/initial conditions, conservation or source terms, and a weak or classical solution class. Moving-domain concentration equations include dilution/transport terms. A stationary law is used only for a time-homogeneous augmented model under existence assumptions; a stationary approximation to a driven system must be justified. Absorbing extinction states are retained. Survival-conditioned, quasi-stationary and unconditional laws are distinct and must not be substituted for one another.

\paragraph{Identification and causal meaning.}
Two parameter values are observationally equivalent when they induce the same full distribution under all allowed observation/intervention designs, not merely the same likelihood value for one sample. A target is identified only if constant on the equivalence class. Gauge symmetries, label switching and reference conventions are quotiented or fixed. Finite-data compatible sets use a declared confidence, likelihood or discrepancy rule; an empty set signals model rejection. Bayesian credible sets are not automatically confidence sets. A prior or regularizer does not create data identification. Bounds intended to resolve a question must be informative at a predeclared width or decision threshold.

A causal contrast uses well-defined intervention versions and a stated causal model. Random assignment, or explicit consistency, overlap and exchangeability assumptions, is required for causal interpretation from observations. Interference and time-varying confounding are represented when present. Fixing a mediator by intervention differs from conditioning on its observed value. Marginal potential-outcome laws do not determine individual cross-world effects or an arbitrary joint coupling. Ethically admissible designs, rather than unconstrained interventions, define the experimental menu.

\paragraph{Metrics, units, optimization and limits.}
Specify the sampling or reference measure for densities and integrals, the log base, and dimensionless arguments of logarithms. Total variation is $\sup_A|P(A)-Q(A)|$; Wasserstein distances require the declared ground metric and appropriate moment conditions; KL divergence uses its stated absolute-continuity convention. All normed readouts are scaled consistently, and error probabilities use $0<\alpha<1$. A ratio has a positive denominator or an explicit undefined/zero convention. Thermodynamic entropy production is nonnegative in expectation under the relevant hypotheses; individual stochastic realizations need not be. An infimum is not an attained minimum without compactness/coercivity and continuity or another existence argument. Fix the order of large-population, long-time and vanishing-mutation limits. Stability claims declare perturbation space and symmetries.

\subsection*{How the provisional difficulty order was assigned}
Four editorial components are scored from 1 (lower burden) to 5 (higher burden): (a) formal or inferential difficulty, (b) access to evidence and identification burden, (c) mechanistic/computational integration, and (d) independent validation and transfer. Their unweighted total orders the entries. For selected entries in the highest bands, an explicit editorial priority breaks equal totals; all other ties are broken by original problem number. Equal totals are \emph{not} evidence of reliably different true difficulties. Scores are judgments, not measurements, completion-time estimates or consensus rankings.

The object ranked is an informative, nondegenerate resolution of the \emph{scoped} scientific target, with the stated budget and validation domain. A sharp useful impossibility result can qualify; an already known identity, vacuous bound, arbitrary favorable model or lookup table cannot. Rank 330 is merely the least difficult relative assessment in this catalogue, not a claim of an easy biological solution. Broad historical titles do not automatically outrank precisely formulated mechanistic tasks.



\end{document}
